\editionPage[\pdfbookmark[0]{Corrigés des défis}{bm-defi-cor-1-parent}]{defi-cor-1}{D01 · Corrigé}{Cent termes, sans cent calculs}{\hyperref[defi-ex-1]{Énoncé p.\,\pageref*{defi-ex-1}}}
\begin{blocEdition}{defi-cor-1}
\editionEtape{Regrouper sans changer l’ordre}On écrit $(-1+2)+(-3+4)+\cdots+(-99+100)$. Chaque parenthèse vaut 1. Les 100 termes forment 50 paires, donc $\res{S=50}$.
\editionEtape{Retirer le dernier terme}S’arrêter à $-99$ revient à enlever le terme $+100$ de la somme précédente : $50-100=\res{-50}$.
\editionEtape{Contrôle}Avec seulement quatre termes, $-1+2-3+4=2$ : on retrouve deux paires. Le raisonnement repose sur la même structure pour les 50 paires, pas sur une supposition tirée du petit exemple.
\end{blocEdition}
\editionPage{defi-cor-2}{D02 · Corrigé}{Retrouver tout le budget}{\hyperref[defi-ex-2]{Énoncé p.\,\pageref*{defi-ex-2}}}
\begin{blocEdition}{defi-cor-2}
\editionEtape{Suivre les restes}Après la première utilisation, il reste $\frac23$ du budget. Après la deuxième, il reste $\frac34\times\frac23=\frac12$ du budget. La troisième conserve $\frac25$ de ce montant, soit $\frac25\times\frac12=\frac15$ du budget initial.
\editionEtape{Revenir au total}Un cinquième du budget vaut 18 euros, donc le budget est $\res{90}$ euros.
\editionEtape{Calculer les parts}Le tiers initial vaut 30 euros. Il reste 60 euros ; leur quart vaut 15 euros. Il reste 45 euros ; leurs trois cinquièmes valent 27 euros. Contrôle : $30+15+27+18=90$. Les fractions successives ne portaient pas sur le même total.
\end{blocEdition}
\editionPage{defi-cor-3}{D03 · Corrigé}{Jusqu’où peut-on doubler ?}{\hyperref[defi-ex-3]{Énoncé p.\,\pageref*{defi-ex-3}}}
\begin{blocEdition}{defi-cor-3}
\editionEtape{Modéliser}Au départ, il y a 1 élément. Après $n$ étapes, il y en a $2^n$ : chaque étape ajoute un facteur 2.
\editionEtape{Encadrer la capacité}$2^{10}=1\,024$, donc $2^{20}=2^{10}\times2^{10}=1\,048\,576$ et $2^{19}=524\,288$. Ainsi $\res{2^{19}<1\,000\,000<2^{20}}$.
\editionEtape{Conclure}On peut effectuer $\res{19}$ étapes complètes ; une vingtième dépasserait la capacité. La valeur de départ ne compte pas comme une étape de doublement.
\end{blocEdition}
\editionPage{defi-cor-4}{D04 · Corrigé}{Racine de deux, un nombre non décimal}{\hyperref[defi-ex-4]{Énoncé p.\,\pageref*{defi-ex-4}}}
\begin{blocEdition}{defi-cor-4}
\editionEtape{1 · Pas un entier}$1^2=1<2<4=2^2$, donc $1<\sqrt2<2$ : aucun entier ne convient.
\editionEtape{2 · Le dernier chiffre}Supposons que $\sqrt2=d$, avec $k\ge1$ chiffres après la virgule, le dernier étant $c$, entre 1 et 9. Alors $d^2$ a $2k$ chiffres après la virgule, et son dernier chiffre est celui de $c^2$. Or les carrés de 1 à 9 se terminent par 1, 4, 9, 6, 5, 6, 9, 4 et 1 : $\res{\text{jamais par 0}}$. Le nombre $d^2$ n’est donc pas entier.
\editionEtape{3 · La contradiction}Pourtant $d^2=2$ est un entier : c’est absurde. L’hypothèse est fausse : $\res{\sqrt2\text{ n’est pas décimal}}$. Par exemple, $1{,}41^2=1{,}9881$ et $1{,}42^2=2{,}0164$ encadrent 2 sans l’atteindre.\label{fig-defi-cor-4}\figureGeo[corrige]{defi-racine}

\end{blocEdition}
\editionPage{defi-cor-5}{D05 · Corrigé}{Comparer sans choisir de valeur}{\hyperref[defi-ex-5]{Énoncé p.\,\pageref*{defi-ex-5}}}
\begin{blocEdition}{defi-cor-5}
\editionEtape{Écrire les aires}$A=3(x+4)=3x+12$ et $B=2(x+6)=2x+12$, en cm$^2$.
\editionEtape{Comparer en général}$A-B=(3x+12)-(2x+12)=x$. Puisque $x>0$, cette différence est positive : $\res{A>B}$ pour toutes les longueurs admises.
\editionEtape{Étudier l’égalité}L’égalité des aires imposerait $x=0$, exclu par l’énoncé. Tester plusieurs nombres n’aurait pas prouvé la conclusion pour tous les nombres positifs ; la différence littérale le permet.\label{fig-defi-cor-5}\figureGeo[corrige]{defi-litteral}

\end{blocEdition}
\editionPage{defi-cor-6}{D06 · Corrigé}{Compter les pièces}{\hyperref[defi-ex-6]{Énoncé p.\,\pageref*{defi-ex-6}}}
\begin{blocEdition}{defi-cor-6}
\editionEtape{Choisir une inconnue}Si $x$ est le nombre de pièces de 2 euros, il y a $42-x$ pièces de 1 euro. Leur valeur totale vaut $2x+(42-x)$.
\editionEtape{Résoudre}L’équation $2x+42-x=68$ se réduit en $x+42=68$, donc $\res{x=26}$. Il y a $\res{16}$ pièces de 1 euro.
\editionEtape{Vérifier le contexte}Les nombres sont entiers et positifs. On compte $26+16=42$ pièces et $2\times26+16=68$ euros. Contrôler seulement la somme d’argent n’aurait pas vérifié le nombre de pièces.
\end{blocEdition}
\editionPage{defi-cor-7}{D07 · Corrigé}{Partager un bénéfice}{\hyperref[defi-ex-7]{Énoncé p.\,\pageref*{defi-ex-7}}}
\begin{blocEdition}{defi-cor-7}
\editionEtape{1 · Première année}$12\,000:18\,000:30\,000=2:3:5$. Une part vaut $9\,000:10=900$ € : les parts sont $\res{1\,800}$ €, $\res{2\,700}$ € et $\res{4\,500}$ €.
\editionEtape{2 · Deuxième année}$12\,000:18\,000:20\,000=6:9:10$, soit 25 parts. Une part vaut $10\,000:25=400$ € : $\res{2\,400}$ €, $\res{3\,600}$ € et $\res{4\,000}$ €.
\editionEtape{3 · Comparer}Les variations sont $+600$ €, $+900$ € et $-500$ € : le $\res{\text{deuxième associé}}$ gagne le plus. Le troisième reçoit moins malgré un bénéfice total plus grand, car sa part du capital a diminué.
\end{blocEdition}
\editionPage{defi-cor-8}{D08 · Corrigé}{Annuler une remise}{\hyperref[defi-ex-8]{Énoncé p.\,\pageref*{defi-ex-8}}}
\begin{blocEdition}{defi-cor-8}
\editionEtape{Coefficients}Baisser de 20\,\% revient à multiplier par 0,8. Pour revenir au prix $P$, il faut multiplier le nouveau prix par un nombre $k$ tel que $0{,}8\times k=1$, soit $k=\frac1{0{,}8}=1{,}25$ : c’est une hausse de $\res{25\,\%}$.
\editionEtape{Généraliser}Pour tout prix $P>0$ : $P\times0{,}8\times1{,}25=P\times1=P$. Le résultat ne dépend pas de $P$.
\editionEtape{Éviter le piège}Une hausse de 20\,\% donnerait $P\times0{,}8\times1{,}2=0{,}96P$ : il manquerait 4\,\% du prix. Les deux pourcentages ne portent pas sur la même base.
\end{blocEdition}
\editionPage{defi-cor-9}{D09 · Corrigé}{Un découpage variable}{\hyperref[defi-ex-9]{Énoncé p.\,\pageref*{defi-ex-9}}}
\begin{blocEdition}{defi-cor-9}
\editionEtape{Établir la formule}L’aire initiale est 180 cm$^2$, donc $\res{A(x)=180-4x^2}$.
\editionEtape{Calculer les trois cas}$A(2)=164$, $A(3)=144$ et $A(4)=\res{116}$ cm$^2$. Le choix est $x=4$ cm.
\editionEtape{Retrouver directement}Retirer $180-116=64$ cm$^2$ en quatre carrés identiques donne 16 cm$^2$ par carré. Le côté positif vaut $\sqrt{16}=4$ cm, et la contrainte $0<x<5$ est respectée.\label{fig-defi-cor-9}\figureGeo[corrige]{coins-retires}

\end{blocEdition}
\editionPage{defi-cor-10}{D10 · Corrigé}{Même moyenne, autre médiane}{\hyperref[defi-ex-10]{Énoncé p.\,\pageref*{defi-ex-10}}}
\begin{blocEdition}{defi-cor-10}
\editionEtape{Moyennes}Les sommes valent 60 et les effectifs 6 : les deux moyennes valent $\res{10}$.
\editionEtape{Médianes et étendues}Pour A, les valeurs centrales sont 4 et 16 : toute valeur entre 4 et 16 est une médiane, par exemple 10 ; l’étendue vaut 12. Pour B, les valeurs centrales valent 8 : la médiane est $\res8$ et l’étendue vaut 28.
\editionEtape{Interpréter}Une même moyenne n’impose ni la même médiane ni la même dispersion. La valeur 28 compense les petites valeurs de B dans la somme. Pour décrire une série, on précise l’indicateur utilisé sans lui faire dire davantage que ce qu’il mesure.
\end{blocEdition}
\editionPage{defi-cor-11}{D11 · Corrigé}{Un jeu de dés équitable ?}{\hyperref[defi-ex-11]{Énoncé p.\,\pageref*{defi-ex-11}}}
\begin{blocEdition}{defi-cor-11}
\editionEtape{1 · Produit pair}Les 36 issues sont équiprobables. Le produit est impair seulement si les deux nombres sont impairs : $3\times3=9$ issues. Donc $P(\text{Léa gagne})=1-\frac9{36}=\res{\frac34}$ : le jeu $\res{\text{n’est pas équitable}}$, il avantage Léa.
\editionEtape{2 · Somme paire}La somme est paire si les deux nombres sont pairs ($9$ issues) ou tous deux impairs ($9$ issues) : $P=\frac{18}{36}=\res{\frac12}$. Le jeu devient $\res{\text{équitable}}$.
\editionEtape{Contrôle}Le tableau corrigé colore les issues gagnantes pour Léa avec la première règle : on en compte 27, soit bien $\frac{27}{36}=\frac34$.\label{fig-defi-cor-11}\figureGeo[corrige]{defi-des}

\end{blocEdition}
\editionPage{defi-cor-12}{D12 · Corrigé}{Six triangles de même aire}{\hyperref[defi-ex-12]{Énoncé p.\,\pageref*{defi-ex-12}}}
\begin{blocEdition}{defi-cor-12}
\editionEtape{Des paires égales}Dans le triangle GBC, $[GI]$ est une médiane : GBI et GIC ont la même aire $a$. De même, GCJ et GJA ont la même aire $b$, et GAK et GKB la même aire $c$.
\editionEtape{Utiliser les grandes médianes}La médiane $[AI]$ partage ABC en ABI et AIC, de même aire : $c+c+a=a+b+b$, donc $c=b$. La médiane $[BJ]$ donne de même $a+a+b=b+c+c$, donc $a=c$.
\editionEtape{Conclure}Ainsi $\res{a=b=c}$ : les six petits triangles ont la même aire, égale au sixième de celle de ABC. Seule la propriété de la médiane, vue en Cinquième, a été utilisée.\label{fig-defi-cor-12}\figureGeo[corrige]{defi-medianes}

\end{blocEdition}
\editionPage{defi-cor-13}{D13 · Corrigé}{Le pont le plus court}{\hyperref[defi-ex-13]{Énoncé p.\,\pageref*{defi-ex-13}}}
\begin{blocEdition}{defi-cor-13}
\editionEtape{Utiliser une translation}La translation qui transforme N en M transforme B en $B'(6\,;4)$. Le quadrilatère MNBB’ est un parallélogramme, donc $NB=MB'$. Le trajet vaut $AM+MB'+1$, car $MN=1$ km.
\editionEtape{Le plus court chemin}$AM+MB'$ est minimal quand M est sur le segment $[AB']$. La droite $(AB')$ passe par les points $(3\,;2)$ et $(6\,;4)$ : M est le point $\res{(3\,;2)}$ de la rive.
\editionEtape{La longueur}$AB'=\sqrt{6^2+4^2}=\sqrt{52}$ km, d’après le théorème de Pythagore. Le trajet minimal vaut $\res{\sqrt{52}+1\approx8{,}21\ \mathrm{km}}$. La translation a transformé le problème en une ligne droite.\label{fig-defi-cor-13}\figureGeo[corrige]{defi-pont}

\end{blocEdition}
\editionPage{defi-cor-14}{D14 · Corrigé}{Varignon et les aires}{\hyperref[defi-ex-14]{Énoncé p.\,\pageref*{defi-ex-14}}}
\begin{blocEdition}{defi-cor-14}
\editionEtape{Un coin}Dans ABD, $[DI]$ est une médiane : $\mathcal A(AID)=\frac12\mathcal A(ABD)$. Dans AID, $[IL]$ est une médiane : $\mathcal A(AIL)=\frac12\mathcal A(AID)=\frac14\mathcal A(ABD)$. De même, $\mathcal A(CKJ)=\frac14\mathcal A(CDB)$.
\editionEtape{Les quatre coins}Les coins A et C totalisent $\frac14\left[\mathcal A(ABD)+\mathcal A(CDB)\right]=\frac14\mathcal A(ABCD)$. De même, les coins B et D totalisent $\frac14\mathcal A(ABCD)$.
\editionEtape{Conclure}Le parallélogramme est ce qui reste : $\mathcal A(IJKL)=\mathcal A(ABCD)-\frac12\mathcal A(ABCD)=\res{\frac12\mathcal A(ABCD)}$. Le résultat ne dépend pas de la forme du quadrilatère.\label{fig-defi-cor-14}\figureGeo[corrige]{defi-varignon}

\end{blocEdition}
\editionPage{defi-cor-15}{D15 · Corrigé}{Le câble et ses attaches}{\hyperref[defi-ex-15]{Énoncé p.\,\pageref*{defi-ex-15}}}
\begin{blocEdition}{defi-cor-15}
\editionEtape{Distance à franchir}D’après le théorème de Pythagore, $d^2=9^2+12^2=225$, donc $d=15$ m.
\editionEtape{Ajouter les réserves}La longueur totale est $15+2\times0{,}75=\res{16{,}5\ \mathrm m}$. La réserve n’est pas un côté du triangle et n’entre pas sous la racine.
\editionEtape{Calculer le coût}$16{,}5\times4{,}80=\res{79{,}20}$ euros. On suppose que le câble peut être acheté au dixième de mètre, sans supplément.\label{fig-defi-cor-15}\figureGeo[corrige]{defi-pyth}

\end{blocEdition}
\editionPage{defi-cor-16}{D16 · Corrigé}{Le plus grand triangle du demi-cercle}{\hyperref[defi-ex-16]{Énoncé p.\,\pageref*{defi-ex-16}}}
\begin{blocEdition}{defi-cor-16}
\editionEtape{1 · Une somme constante}M est sur le cercle de diamètre $[AB]$ : le triangle AMB est rectangle en M. Le théorème de Pythagore donne $\res{AM^2+BM^2=AB^2=100}$.
\editionEtape{2 · Majorer l’aire}Soit H le pied de la hauteur issue de M. Dans le triangle MHO, rectangle en H, l’hypoténuse $[OM]$ est le plus long côté : $MH\le OM=5$ cm. L’aire vaut $\frac{AB\times MH}2\le\frac{10\times5}2=\res{25\ \mathrm{cm}^2}$.
\editionEtape{Le cas d’égalité}L’aire vaut 25 cm$^2$ lorsque $H=O$, c’est-à-dire lorsque $(OM)\perp(AB)$ : M est alors à égale distance de A et B, et $AM=BM=\sqrt{50}$ cm.\label{fig-defi-cor-16}\figureGeo[corrige]{defi-cercle}

\end{blocEdition}
\editionPage{defi-cor-17}{D17 · Corrigé}{Une pyramide de verre}{\hyperref[defi-ex-17]{Énoncé p.\,\pageref*{defi-ex-17}}}
\begin{blocEdition}{defi-cor-17}
\editionEtape{1 · Volume}$V=\frac{35^2\times21}3=\frac{1\,225\times21}3=\res{8\,575\ \mathrm m^3}$.
\editionEtape{2 · Hauteur d’une face}Dans le triangle SOM rectangle en O : $SO=21$ m et $OM=17{,}5$ m. Pythagore donne $SM^2=441+306{,}25=747{,}25$, donc $SM\approx\res{27{,}34\ \mathrm m}$.
\editionEtape{3 · Surface vitrée}Chaque face a pour aire $\frac{35\times SM}2$. Les quatre faces totalisent $70\times SM\approx\res{1\,914\ \mathrm m^2}$, avec la valeur non arrondie de SM. Le modèle ignore l’épaisseur des montants.\label{fig-defi-cor-17}\figureGeo[corrige]{defi-pyramide}

\end{blocEdition}
\editionPage{defi-cor-18}{D18 · Corrigé}{La suite de Syracuse}{\hyperref[defi-ex-18]{Énoncé p.\,\pageref*{defi-ex-18}}}
\begin{blocEdition}{defi-cor-18}
\editionEtape{1 · Les valeurs}Pour 7 : $\res{22,\ 11,\ 34,\ 17,\ 52,\ 26,\ 13,\ 40,\ 20,\ 10,\ 5,\ 16,\ 8,\ 4,\ 2,\ 1}$. Chaque valeur ne dépend que de la précédente.
\editionEtape{2 · Le retour à 1}On obtient 1 au $\res{16^\text{e}}$ tour. Ensuite, on obtiendrait 4, 2, 1, puis de nouveau 4, 2, 1 : la suite devient périodique. Personne n’a encore démontré qu’on atteint 1 pour tout entier de départ.
\editionEtape{3 · Une seule variable}Chaque instruction remplace n par une nouvelle valeur calculée à partir de l’ancienne ; on n’a pas besoin de garder les valeurs précédentes. En Troisième, une boucle « répéter jusqu’à n = 1 » permettra de s’arrêter au bon moment.\label{fig-defi-cor-18}\figureGeo[corrige]{defi-script}

\end{blocEdition}
